\begin{lemma}[Situation 7.2, Subcase 3.1]
Every \noderef{HierarchySituationTwoCaseThreeOneData} is impossible.
Equivalently, normalized Case 3 cannot occur when the last-edge color
\(\theta=\varphi_n(e_p)\) is present at \(y_{p-1}\).
\end{lemma}
\begin{proof}
The preceding \noderef{HierarchySituationTwoCaseThreeNormalization}
is the branch-neutral source of condition~(20); here that condition is
part of the supplied datum and we prove the subsequent `.1' branch.
Write \(H=C.\mathit{hierarchy}\), let \(q\) be its last stored index, and
put
\[
 A=\noderef{HierarchyStarredStage}(H,q),\qquad
 \ell=\noderef{HierarchyFullTailLength}(C),
\]
\[
 B=\noderef{HierarchyFullTailPrefix}(C,\ell-2),\qquad
 B^+=\noderef{HierarchyFullTailPrefix}(C,\ell-1).
                                                               \tag{1}
\]
The exact position fields of the datum and \(\ell\geq2\) give
\[
 A\preccurlyeq B\prec B^+\prec T,\qquad
 |B|=|T|-2,\quad |B^+|=|T|-1,                                \tag{2}
\]
as entrywise \noderef{TreeSequenceInitialSegment}s.  In particular the
two final displayed vertices are \(y_{p-1}\) and \(y_p\), and the final
individual edge is \(e_p\).  Apply
\noderef{HierarchyCounterexamplePenultimateElementary} to the properness,
exact one-uncoloredness, complete lower-rung package, fewer-divider
induction, lexicographic minimality, and large-set fields of the supplied
\noderef{HierarchyMinimumCounterexampleContext}.  Its conclusion is
\[
             V(B^+)\text{ is elementary under }\varphi_n.    \tag{3}
\]

We first force the exception in manuscript (23).  Form a local
\noderef{HierarchySituationTwoTailData} by copying from \(D\) the lower
bound on \(\ell\), the full path-number equation, the three exact final
indices and their position equations, the non-incidence of \(e_p\) with
\(y_{p-1}\), and \(\varphi_n(e_p)=\theta\), and by using precisely
\(D.\mathit{thetaPresentPenultimate}\) for its last field.  This is a
fieldwise construction, not an identification of the common-tail and
`.1' data structures.  Invoke
\noderef{HierarchySituationTwoReorderedTreesCertified} on this adapter.
Retain its exact reordered tree \(T^-\), the PE or exact-seed SE/RE
augmentation route, the actual
\noderef{CorrespondingExtendedTashkinovTree}, maximum property, rebuilt
base \(\mathcal B^-\), and hierarchy \(H^-\).  The returned
\noderef{HierarchySituationTwoReorderedTrees} says
\[
       T^-=B+(e_p,y_p),\qquad
       V(T^-)=V(T)\setminus\{y_{p-1}\},                     \tag{4}
\]
and the certificate says that \(H^-\) has the old core, PE closure when
applicable, divider count, all ordinary stages, and all numerical divider
positions.

Suppose, for a contradiction, that there is no \(\eta\in D_{n,q}\)
which is missing at \(y_{p-1}\) and satisfies
\(\theta\in\Gamma^q_\eta\), where
\[
 D_{n,q}=\noderef{HierarchyStoredColorsAt}(H,q).
                                                               \tag{5}
\]
We construct a \noderef{GoodHierarchy} on \(H^-\).  Transport the old
Gamma function across the returned equality of divider counts.  Exact
agreement of the raw stored arrays, the base core, and the retained
ordinary stages identifies every computed set (5) and every
\noderef{HierarchyStageMissingColors}.  It therefore transports the
two-color cardinality, pairwise disjointness, unique refresh,
parallel-base, and minus-closed fields.  For avoidance at an index below
\(q\), both endpoints of the interval in
\noderef{HierarchySegmentColors} are retained, so the old avoidance field
applies literally.

For avoidance at \(q\), fix \(\zeta\in D_{n,q}\) and
\(c\in\Gamma^q_\zeta\).  If \(\zeta\) has a missing witness in \(B\),
its least missing position and all earlier exact edges occur in the
unchanged prefix \(B\), so old avoidance applies.  Otherwise the only
edge moved to an earlier position by (4) is \(e_p\), whose old color is
\(\theta\).  Hence only \(c=\theta\) can create a new occurrence before
the recalculated cutoff.  If \(c=\theta\), the assumed failure of the
exception says that \(\zeta\) is not missing at \(y_{p-1}\).  Old
avoidance of the old last edge then forces the first \(\zeta\)-missing
position to occur no later than \(B\); otherwise \(e_p\), which precedes
either the witness \(y_p\) or the terminal fallback, would already be a
forbidden \(\theta\)-edge.  This contradicts the present no-witness
case.  Thus terminal avoidance also holds and all fields of the claimed
good hierarchy on \(H^-\) have been checked.

Condition (20a) supplies a vertex of \(B\) missing \(\alpha\), while the
raw Case 3 field supplies the distinct terminal witness \(y_p\).  Both
belong to \(T^-\) by (2)--(4).  They make \(V(T^-)\) non-elementary.
Together with the returned ETT, base, hierarchy, goodness, and maximum
property they form a \noderef{HierarchyCounterexample} \(C^-\).  It has
the old rung and divider count.  Its last starred stage is \(A\), while
\(|T^-|=|T|-1\), so its full-tail length is \(\ell-1\).  Expanding
\noderef{HierarchyPathNumber}, its value is at most that full-tail
length: a nonempty candidate minimum lies in its defining interval, and
the empty-candidate fallback is exactly the full-tail length.  Since the
datum says that the old path number is \(\ell\), the first coordinate of
\noderef{HierarchyCounterexampleMeasure}(C^-)\) is strictly smaller.
Thus \noderef{HierarchyCounterexampleLess}(C^-,C)\), contradicting the
old lexicographic-minimality field.  We have proved that there exists a
color \(\eta_m\) such that
\[
 \eta_m\in D_{n,q},\qquad
 \eta_m\in\overline\varphi_n(y_{p-1}),\qquad
 \theta\in\Gamma^q_{\eta_m}.                                \tag{6}
\]

Let \(v_\alpha\in V(B)\) be a vertex missing \(\alpha\).  Elementarity
(3) and (6) imply \(\alpha\ne\eta_m\), and also imply that no vertex of
\(B\) misses \(\eta_m\).  Membership in (5) says that \(\eta_m\) is not
missing on \(A\), so its totalized
\noderef{HierarchyFirstMissingPosition} is exactly \(y_{p-1}\).  The
avoidance field of \noderef{GoodHierarchy}, applied to (6), now gives
both that \(\theta\) is in the appropriate stage-missing set and
\[
       \theta\notin\varphi_n\langle B^+-A\rangle .           \tag{7}
\]
This is equation (24), with the interval interpreted through the exact
edge identities in \noderef{HierarchySegmentColors}.

We perform the first change.  The colors \(\theta\) and \(\alpha\) are
distinct: \(e_p\) is incident with \(y_p\), has color \(\theta\), and
\(\alpha\) is missing there.  Both colors are missing on \(B^+\):
\(\alpha\) has its witness in \(B\), while the stage-missing conclusion
above gives a \(\theta\)-witness in the last ordinary stage, hence in
\(A\).  Apply \noderef{ChangeLemma} with ordered colors
\((\theta,\alpha)\).  When its conditional order premise is active,
\(\theta\)'s first witness is in \(A\), whereas \(\alpha\)'s is in
\(B\setminus A\), and (7) excludes \(\theta\) from the whole required
interval.  For \(q>0\), \(\theta\) is missing on \(T_{n,q}\); for
\(q=0\), it is missing on \(T_n\).  These are exactly the two Cases 1/2
alternatives.  Using (3) and all fields of the minimum context, the lemma
gives
\[
 K_{\alpha,\theta}(v_\alpha)
       =K_{\alpha,\theta}(v_\theta),                         \tag{8}
\]
where \(v_\theta\) is the first \(\theta\)-missing vertex.

The component in (8) is disjoint from
\(Q_1=K_{\alpha,\theta}(y_p)\).  If
\(v_\alpha=v_\theta\), both colors are missing at the common vertex;
unfolding \noderef{MissingColors} and
\noderef{TwoColorIncidenceGraph} makes its component an isolated vertex,
which cannot contain the distinct \(y_p\).  Otherwise
\(v_\theta,v_\alpha,y_p\) are three distinct vertices of affected-color
incidence degree at most one, missing respectively
\(\theta,\alpha,\alpha\).  A common component would contradict the
finite spanning path in \noderef{TwoColorComponentSpanningPath} and the
degree bound in \noderef{TwoColorIncidenceDegreeAtMostTwo}.  Thus
\noderef{TwoColorComponentsDisjoint} gives disjointness in both cases,
and the spanning-path result supplies a possibly singleton
\noderef{TwoColorPathComponent} for \(Q_1\).

Define
\[
 \mu_1=\noderef{TwoColorComponentSwap}
       (G,\varphi_n,\theta,\alpha,\operatorname{inl}y_p).    \tag{9}
\]
Before using any stability consequence, apply
\noderef{StableChangeCasesOneTwo} with ordered colors
\((\theta,\alpha)\), (3), (7)--(9), the just-proved disjointness, the
stage-missing alternative for \(\theta\), and the six minimum-context
inputs.  Retain its \noderef{HierarchyWeaklyStable} and
\noderef{ModColoring} conclusions.

We verify every input of
\noderef{StableChangeCasesOneTwoSplitTail} with split \(X=B^+\).  First,
condition (20a) excludes \(\alpha\) from every exact edge of \(T-A\),
and (7) excludes \(\theta\) from every exact edge of \(B^+-A\).
Consequently
\[
 \mu_1(T.\mathit{edge}(r))=\varphi_n(T.\mathit{edge}(r))
       \quad(|A|\le r<|B^+|).                               \tag{10}
\]
Second, missing sets agree at every vertex of \(B^+\).  Outside \(Q_1\)
no incident edge changes.  A vertex of \(Q_1\cap V(B^+)\) cannot be an
endpoint: by (3), an endpoint missing either affected color would be the
corresponding witness in the disjoint component (8).  It is therefore
internal, with exactly one incident edge of each affected color by
properness, and (9) exchanges those two edges.  All other colors are
fixed.  Hence
\[
 \overline\mu_1(u)=\overline\varphi_n(u)
                 \quad(u\in V(B^+)).                        \tag{11}
\]
In particular both affected colors remain missing on \(B^+\), at the
witnesses belonging to (8).

We next produce the exact route through \(B^+\).  Inspect the final
certificate in the actual \noderef{TashkinovSeries}.  Its PE constructor
supplies a \(\varphi_n\)-\noderef{TashkinovClosure} from \(T_n\); its
SE/RE constructors supply the individual edge \(f_n\), the literal
singleton connecting-edge equality, the exact
\noderef{TreeSequenceAdjoinEdge} seed \(T_n+f_n\), and a closure from
that seed.  In PE, the base-stage chain and, at divider zero, the left
segment of \noderef{OrderedTreeUnion}, put \(T_n\) below \(A\).  In
SE/RE with \(q=0\), the nonparallel base gives \(A=T_n\), and
\(\ell\ge2\) puts the one-edge seed below \(B^+\).  In SE/RE with
\(q>0\), \noderef{HierarchyNonparallelFirstEdge} identifies the exact
seed edge with the first edge of hierarchy stage one; comparing the two
prefixes of the common final series tree puts
\(T_n+f_n\preccurlyeq A\preccurlyeq B^+\).  In every case,
\noderef{TashkinovAugmentationSequenceRestrict} therefore restricts the
old closure route to a route from the eligible PE source or exact SE/RE
seed through \(B^+\).

Transport that prefix route to \(\mu_1\) step by step.  Up to \(A\),
unfold the retained \noderef{HierarchyWeaklyStable}: for positive
divider count its \noderef{TreeColoringInvariant} on \(A\) preserves
each predecessor missing witness and every added edge value; at divider
zero its PE combined-stability clauses on \(T_n\) and the PE closure,
together with \noderef{CombinedStableClosurePreserved} and the
\noderef{OrderedTreeUnion} entry decomposition, fix every displayed edge
of \(A\); while in SE/RE the nonparallel equality \(A=T_n\) leaves only
the exact seed, whose incident edge is protected by the terminal
\noderef{StableColoring} clause.  Above \(A\), equations (10)--(11)
verify the edge and missing-witness hypotheses of
\noderef{TashkinovAugmentationSequenceColoringTransport}.  Concatenating
the transported pieces gives exactly the required \(\mu_1\)-route to
\(B^+\), with the same PE source or the same individual SE/RE seed.

The old suffix \(B^+\to T\) is the single
\noderef{TashkinovAugmentation} supplied by the exact final adjoin
certificate retained in (4): its new edge is \(e_p\), its color is
\(\theta\), and \(\theta\) is missing already in \(A\subseteq B^+\).
This is the required original-color suffix.  Under \(\mu_1\), both
affected colors are missing on \(B^+\), as observed after (11), so the
suffix is also transportable by the component-swap theorem used inside
the split-tail provider.

It remains to verify stored positions and the Gamma cutoff.  The color
\(\theta\) is not stored, because (6) and goodness make it missing on
the last ordinary stage.  For \(\alpha\), the first two alternatives of
condition~(20) say directly that it is not stored.  In the third
alternative it is stored, but it has a missing witness in \(B\), and
(11) identifies all missing-position candidates through that witness;
thus its least position is unchanged.  Every other stored color is
unaffected everywhere.  Unfolding
\noderef{HierarchyFirstMissingPosition}, including its terminal fallback,
therefore proves the required equality for every stored color.  Finally,
condition~(20) and the reserved/stored disjointness in
\noderef{GoodHierarchyReservedDisjoint} put \(\alpha\) in no Gamma pair.
Pairwise
disjointness and (6) say that a Gamma pair containing \(\theta\) is the
pair of \(\eta_m\); its first missing position is \(y_{p-1}\), whose
index is \(|B^+|\).  This is exactly the quantified Gamma-cutoff premise.

The split-tail theorem now returns an actual corresponding ETT on the
same exact \(T\), maximum property, a rebuilt base \(\mathcal B_1\), a
hierarchy \(H_1\), a good-hierarchy witness \(g_1\), and
\[
          \noderef{SameGoodHierarchy}(C.\mathit{good},g_1). \tag{12}
\]
The swap makes \(\theta\) missing at \(y_p\): the old missingness of
\(\alpha\), properness, and the incident \(\theta\)-edge \(e_p\) show
that the unique affected incident edge changes from \(\theta\) to
\(\alpha\).  Disjointness preserves the \(\theta\)-missing witness
\(v_\theta\in B^+\).  These distinct witnesses make the returned
hierarchy non-elementary and hence give an actual
\noderef{HierarchyCounterexample} \(C_1\).

Package the returned data as a
\noderef{HierarchyReplacementPresentation} \(P_1\), checking its
minimum context field by field.  Properness follows from
\noderef{TwoColorComponentSwapProper}, and the definition of the swap
fixes `none', so \noderef{IsOneEdgeUncolored} is retained.  Rung
agreement transports the complete universally quantified lower-rung
six-pack.  Equation (12) transports the divider count, so the old
fewer-divider field applies.  The exact terminal tree is still \(T\),
and (12) identifies all stages and divider positions; expanding
\noderef{HierarchyPathNumber} and
\noderef{HierarchyCounterexampleMeasure} therefore gives equality of
both measure coordinates, after which the old lexicographic-minimality
field applies.  Finally, the replacement-stability field in
\noderef{CorrespondingExtendedTashkinovTree}, unfolded through
\noderef{StableColoring} at the terminal ladder index, identifies its
missing-color set and transports the \(2n+11\) estimate.  Thus the
properness, exact uncolored edge, all six lower-rung assertions,
fewer-divider induction, minimality, and large-set bound are all present;
none is copied across a changed type without its equality.

We now make the second change inside \(P_1\).  Equations (11)--(12)
transport (6): \(\eta_m\) is stored, is missing at \(y_{p-1}\), and its
Gamma pair contains \(\theta\).  It is not missing on \(A\), and
elementarity of \(B^+\), obtained again from
\(P_1.\mathit{minimumContext}\) by
\noderef{HierarchyCounterexamplePenultimateElementary}, makes
\(y_{p-1}\) its first missing vertex.  The color \(\theta\) is missing
at \(v_\theta\) and \(y_p\).  Moreover, (7), (10), and (11) give
\[
 \alpha,\theta\notin\mu_1\langle B^+-A\rangle,\qquad
 \mu_1(e_p)=\alpha.                                        \tag{13}
\]

Apply \noderef{ChangeLemma} relative to \(P_1\), with ordered colors
\((\theta,\eta_m)\).  The first color is stage-missing; the second has
first missing vertex \(y_{p-1}\); and (13) supplies the exact required
unused interval.  The positive-divider or divider-zero Cases 1/2
alternative again uses \(\theta\), not \(\eta_m\).  The conclusion is
\[
 K^{\mu_1}_{\eta_m,\theta}(v_\theta)
       =K^{\mu_1}_{\eta_m,\theta}(y_{p-1}).                  \tag{14}
\]
This component is disjoint from the one rooted at \(y_p\): the three
vertices \(v_\theta,y_{p-1},y_p\) are distinct by (2), storedness of
\(\eta_m\), and the location of \(v_\theta\) in \(A\); they miss
respectively \(\theta,\eta_m,\theta\), so a common spanning path would
again have three endpoints.  Define
\[
 \mu_2=\noderef{TwoColorComponentSwap}
       (G,\mu_1,\theta,\eta_m,\operatorname{inl}y_p).        \tag{15}
\]
Apply \noderef{StableChangeCasesOneTwo} immediately, relative to the
actual presentation \(P_1\), before protecting any exceptional edge.
Use (14), the component separation, and \(\theta\) as the stage-missing
Cases 1/2 color.  Retain its weak-stability conclusion.

We verify the full-tail transport premises separately.  The same
endpoint/internal-vertex argument as for (11), now using (14), proves
\[
 \overline\mu_2(u)=\overline\mu_1(u)
                    \quad(u\in V(B^+)).                     \tag{16}
\]
The color \(\theta\) occurs on no exact tail edge under \(\mu_1\), by
(13) and the changed value of \(e_p\).  Apply
\noderef{HierarchyTailMissingColorExclusion} to \(\eta_m\), whose first
missing vertex is \(y_{p-1}\).  It excludes \(\eta_m\) from every exact
edge of \(T-A\), except possibly the individual divider-zero SE/RE seed
carrying a raw stored color.  The weak stability was obtained before this
argument; in the exceptional case it unfolds to a
\noderef{StableColoring} on \(T_n\), and the exact seed is a
\noderef{TreeIncidentEdge}, so its controlled incident-edge clause fixes
that individual edge.  The final edge has color \(\alpha\), distinct
from \(\eta_m\) by (3) and (6), and distinct from \(\theta\) because
\(\alpha\) was missing at the endpoint of the old \(\theta\)-edge.
Consequently every individual tail edge is fixed:
\[
 \mu_2(T.\mathit{edge}(r))=\mu_1(T.\mathit{edge}(r))
                    \quad(|A|\le r<|T|).                    \tag{17}
\]

For the stored first-position equality, \(\theta\) is not stored.  The
only affected stored color is \(\eta_m\); its first missing position is
\(y_{p-1}\), and (16) identifies every candidate through that witness,
so its least position remains \(y_{p-1}\).  All other stored colors are
unaffected.  Expanding \noderef{HierarchyFirstMissingPosition}, including
the empty-candidate fallback, proves the full quantified equality.
Equations (16)--(17), this equality, (14)--(15), and the stage-missing
alternative verify every premise of
\noderef{StableChangeCasesOneTwoFullTail}.  It returns an actual
corresponding same-rung maximum ETT on \(T\), a base
\(\mathcal B_2\), hierarchy \(H_2\), goodness \(g_2\), and same-good
data relative to \(P_1\).

Construct a second original-relative replacement presentation \(P_2\).
Compose the two correspondence certificates entry by entry: replacement
ladder trees, raw stored sets, exact connecting-edge sets, and extension
kinds compose by equality, and their positive-stage stable-coloring
relations compose by \noderef{sc2}.  Compose the two same-good relations
field by field after the canonical divider-index casts.  The swap (15)
makes \(\eta_m\) missing at \(y_p\), while (16) retains its missingness
at the distinct \(y_{p-1}\); these witnesses construct the required
counterexample.  Properness and exact one-uncoloredness follow as for
\(P_1\).  Rung agreement transports the complete lower-rung package,
composed same-good transports the fewer-divider field and equality of
both measure coordinates before minimality is applied, and terminal
stable-coloring transitivity transports the large-set bound.  Hence all
six minimum-context components are again supplied explicitly.

Finally form a genuine
\noderef{HierarchySituationTwoCaseOneData} for \(P_2\).  Copy the exact
three final indices, their position equations, and the non-incidence of
the last individual edge.  The terminal tree and all hierarchy stages
are unchanged, so the full-tail length and path-number equation are the
old ones.  Use the original color \(\alpha\) as the new tail color:
(13) and (17) give \(\mu_2(e_p)=\alpha\).  The old \(\alpha\)-missing
witness in \(B\), elementarity (3), and the two missing-set equalities
(11) and (16) show that \(\alpha\) is present at \(y_{p-1}\), as the
new `.1' tail datum requires.  Use \(\eta_m\) as the Case 1 color; it is
missing at both \(y_{p-1}\) and \(y_p\) and remains stored under the
composed same-good relation.

The pair \(g_2.\Gamma^q_{\eta_m}\) has cardinality two by
\noderef{GoodHierarchy} and still contains the old \(\theta\).  Choose
its other member \(\gamma\); then
\[
       g_2.\Gamma^q_{\eta_m}=\{\theta,\gamma\},\qquad
       \theta\ne\gamma.                                    \tag{18}
\]
These are exactly the remaining pair fields of the Case 1 datum.  Apply
only the already proved \noderef{HierarchySituationTwoSubcaseOneOne} to
this actual datum and the actual minimum context of \(P_2\).  It gives
the required contradiction.
\end{proof}
